\documentclass[11pt]{amsart}
\pagestyle{plain}
\raggedbottom
\usepackage[dvipsnames]{xcolor}
\usepackage[utf8]{inputenc}
\usepackage{amssymb,amsmath,amsthm,amsfonts,fullpage,float,latexsym,bbm,microtype,cite,fancyvrb}
\usepackage{enumerate}
\usepackage{hyperref}
\hypersetup{colorlinks=true, linkcolor=blue, citecolor=magenta, filecolor=magenta, urlcolor=magenta}
\usepackage{bm,mathrsfs}

\newtheorem{theorem}{Theorem}
\newtheorem{proposition}[theorem]{Proposition}
\newtheorem{claim}[theorem]{Claim}
\newtheorem{conjecture}[theorem]{Conjecture}
\newtheorem{lemma}[theorem]{Lemma}
\newtheorem{corollary}[theorem]{Corollary}
\theoremstyle{definition}
\newtheorem{remark}[theorem]{Remark}
\newtheorem{definition}[theorem]{Definition}
\newtheorem{observation}[theorem]{Observation}
\newtheorem{example}[theorem]{Example}

\renewcommand{\emptyset}{\varnothing}
\newcommand{\C}{\mathbb{C}}
\newcommand{\Z}{\mathbb{Z}}
\DeclareMathOperator{\Hilb}{Hilb}
\DeclareMathOperator{\OSP}{OSP}
\DeclareMathOperator{\inv}{inv}
\DeclareMathOperator{\Inv}{Inv}
\DeclareMathOperator{\sgn}{sgn}
\DeclareMathOperator{\blocks}{blocks}
\DeclareRobustCommand{\qbinom}{\genfrac[]{0pt}{}}

\begin{document}

\title{A combinatorial interpretation of the hook coefficients of the Hilbert
series of the multigraded coinvariant ring}
\maketitle

\section{The Problem}

Let $m$ and $n$ be positive integers and let $R_n^{(m)}$ be the coinvariant
algebra with $m$ sets of $n$ variables. Its multigraded Hilbert series admits a
Schur expansion
$$
\Hilb(R_n^{(m)};q_1,\ldots ,q_m)= \sum_{\lambda}
c_{\lambda}(n)\,s_\lambda(q_1,\ldots ,q_m),
$$
with nonnegative coefficients $c_\lambda(n)$ that do not depend on $m$. We give a
combinatorial interpretation of $c_\lambda(n)$ when $\lambda$ is a \emph{hook
shape} partition, i.e.\ $\lambda=(a,1^b)$ with $a\ge1,\ b\ge0$, or
$\lambda=\varnothing$.

\medskip
\noindent\textbf{Answer (Theorem~\ref{thm:improved-coefficient-formula}).}
For $a \geq 1$ and $b \geq 0$,
$$
c_{(a,1^b)}(n)
= \#\{w \in \OSP(n,n-b)\ :\ \inv(w) = a,\ w \text{ is of type I} \},
$$
where $\OSP(n,n-b)$ is the set of ordered set partitions of $\{1,\dots,n\}$ into
$n-b$ blocks, $\inv$ is the inversion statistic of Sagan--Swanson, and ``type I''
is the mergeability condition recalled in Section~\ref{sec:phi}. Moreover
$c_\varnothing(n)=1$ and $c_\lambda(n)=0$ for every non-hook $\lambda$.

\section{Background}

\subsection{The two coinvariant rings}
The \textbf{coinvariant ring with $m$ sets of variables} is
\begin{equation}\label{eq:m-coinv}
R_n^{(m)} = \C[x_i^{(j)}:1\le i\le n,\ 1\le j\le m]
/\big\langle \C[x_i^{(j)}]_+^{\mathfrak{S}_n} \big\rangle,
\end{equation}
where the subscript $+$ denotes the polynomials with vanishing constant term and
the superscript $\mathfrak{S}_n$ denotes invariance under the diagonal action of
$\mathfrak{S}_n$ permuting the $n$ variables within each of the $m$ sets. It is
$\Z_{\ge0}^m$-multigraded, the $j$-th degree counting the total degree in the
$j$-th set $x_1^{(j)},\dots,x_n^{(j)}$. The associated multigraded Hilbert series
is
\begin{equation}
\Hilb(R_n^{(m)}; q_1,\ldots,q_m) :=
\sum_{r_1,\ldots,r_m \geq 0}\dim\big((R_n^{(m)})_{r_1,\ldots,r_m}\big)
q_1^{r_1}\cdots q_m^{r_m}.
\end{equation}

The \textbf{superspace coinvariant ring} is
\begin{equation}\label{eq:super-coinv}
R_n^{(1,1)} = \C[x_1,\ldots,x_n, \theta_1, \ldots,\theta_n]
/\big\langle \C[x_1,\ldots,x_n, \theta_1, \ldots,\theta_n]_+^{\mathfrak{S}_n}
\big\rangle,
\end{equation}
where the $x_i$ are commuting and the $\theta_i$ are anticommuting
($\theta_i\theta_j=-\theta_j\theta_i$, so $\theta_i^2=0$), and $\mathfrak{S}_n$
acts diagonally. It is bigraded: each $x_i$ contributes to the $q$-degree and
each $\theta_i$ to the $u$-degree. Its bigraded Hilbert series is
\begin{equation}
\Hilb(R_n^{(1,1)}; q;u) :=
\sum_{r,s\ge0}\dim\big((R_n^{(1,1)})_{r,s}\big)q^{r}u^{s}.
\end{equation}

\subsection{Symmetric-function tools}
Let $[k]_q := 1+q+\cdots+q^{k-1}$, $[k]_q!:=[k]_q[k-1]_q\cdots[1]_q$ with
$[0]_q!:=1$, and define the $q$-Stirling numbers of the second kind $S[n,k]$ by
$S[0,k]=\delta_{0,k}$, $S[n,0]=\delta_{n,0}$, and
\begin{equation}\label{eq:stirling}
S[n,k] = S[n-1,k-1] + [k]_q\,S[n-1,k]\qquad (0\le k\le n).
\end{equation}
We write $s_\lambda$ for a Schur function and $s_{\mu/\nu}$ for a skew Schur
function. The (single super-alphabet) \textbf{super Schur function} in one
commuting variable $q$ and one anticommuting variable $u$ is
\begin{equation}\label{eq:superschur-def}
s_\lambda(q/u) = \sum_{\nu \subseteq \lambda}
s_\nu(q)\,s_{\lambda'/\nu'}(u),
\end{equation}
where $\lambda'$ is the conjugate (transpose) partition and the sum is over
partitions $\nu$ contained in $\lambda$
(see \cite[Sections A.2.2 and 3.2]{ChengWangBook}).

\subsection{Ordered set partitions and the inversion statistic}
An \textbf{ordered set partition} of $[n]:=\{1,\dots,n\}$ into $k$ blocks is a
sequence $w=(S_1/S_2/\cdots/S_k)$ of pairwise disjoint nonempty sets with
$S_1\cup\cdots\cup S_k=[n]$. We write $\OSP(n,k)$ for the set of these, and
$\blocks(w):=k$. Following Sagan--Swanson \cite[Section 5.1]{SaganSwanson2024},
the inversion set and statistic are
\begin{equation}
\Inv(w) = \{(s,S_j)\ :\ s \in S_i \text{ for some } i < j
\text{ and } s \geq \min S_j\},\qquad \inv(w):=\#\Inv(w).
\end{equation}

\section{Reduction to the superspace coinvariant ring}\label{sec:reduction}

This section justifies the reduction of the hook-coefficient problem to a
bigraded computation in $R_n^{(1,1)}$, by invoking two external theorems and
verifying their hypotheses.

\begin{proposition}[Bergeron]\label{prop:bergeron}
Fix $n\ge1$. The multigraded Hilbert series of $R_n^{(m)}$ is a symmetric
polynomial in $q_1,\dots,q_m$, hence has a Schur expansion
\begin{equation}\label{eq:bergeron-expansion}
\Hilb(R_n^{(m)}; q_1,\ldots,q_m)
= \sum_{\lambda:\ \ell(\lambda)\le m} c_\lambda(n)\, s_\lambda(q_1,\ldots,q_m),
\end{equation}
with nonnegative integer coefficients $c_\lambda(n)$ that are independent of $m$.
Moreover every coefficient $c_\lambda(n)$ that can appear for some $m$ already
appears once $m\ge n$.
\end{proposition}

\begin{proof}
This is the content of \cite{Bergeron2020}. The diagonal $\mathfrak{S}_n$-action
permutes the $m$ sets of variables symmetrically; the $\GL_m$ (equivalently
$\mathfrak{S}_m$-symmetric) structure on the polarized coinvariants forces the
multigraded Hilbert series to be a symmetric polynomial in $q_1,\dots,q_m$ whose
Schur expansion has coefficients $c_\lambda(n)\ge0$ independent of $m$. A Schur
polynomial $s_\lambda(q_1,\dots,q_m)$ is nonzero precisely when
$\ell(\lambda)\le m$, which both restricts the sum in
\eqref{eq:bergeron-expansion} to $\ell(\lambda)\le m$ and shows that the
coefficient of every $\lambda$ with $\ell(\lambda)\le n$ is realized once
$m\ge n$. Since $R_n^{(m)}$ is a quotient of a polynomial ring by an ideal
generated in positive degree, all $\dim$'s are finite in each multidegree, so the
$c_\lambda(n)$ are nonnegative integers.
\end{proof}

Because the $c_\lambda(n)$ are independent of $m$ and all are visible at $m=n$,
it suffices to determine $c_\lambda(n)$ for each fixed $n$, e.g.\ from the case
$m=n$. We now isolate the hook coefficients through diagonal supersymmetry. Let
$P(1,1,n)$ denote the set of partitions $\lambda$ with $\lambda_2\le1$ and
$\ell(\lambda)\le n$.

\begin{lemma}\label{lem:P11hooks}
$P(1,1,n)$ is exactly the set of hook partitions of length at most $n$, i.e.\
$P(1,1,n)=\{\varnothing\}\cup\{(a,1^b): a\ge1,\ b\ge0,\ b+1\le n\}$.
\end{lemma}

\begin{proof}
A partition $\lambda$ has $\lambda_2\le1$ iff all parts after the first are $\le1$,
i.e.\ $\lambda=\varnothing$, or $\lambda=(a,1^b)$ with $a\ge1$ and $b\ge0$; these
are precisely the hook shapes. The extra constraint $\ell(\lambda)\le n$ is
$b+1\le n$ for $(a,1^b)$ (and is vacuous for $\varnothing$).
\end{proof}

\begin{proposition}[Diagonal supersymmetry; Lentfer]\label{prop:ds}
Fix $n\ge1$. With the same universal coefficients $c_\lambda(n)$ as in
Proposition~\ref{prop:bergeron},
\begin{equation}\label{eq:coeffs-Hilbert}
\Hilb(R_n^{(1,1)}; q;u) = \sum_{\lambda \in P(1,1,n)}
c_{\lambda}(n)\, s_\lambda(q/u).
\end{equation}
\end{proposition}

\begin{proof}
This is the specialization to $R_n^{(1,1)}$ of the diagonal supersymmetry result,
conjectured in \cite{Bergeron2020} and proved in \cite{Lentfer-Supersymmetry}.
The statement is that the same universal coefficients $c_\lambda(n)$ govern the
bigraded super-coinvariant Hilbert series, where each ordinary Schur function
$s_\lambda$ is replaced by the super Schur function $s_\lambda(q/u)$, and only the
hook indices $\lambda\in P(1,1,n)$ survive (by Lemma~\ref{lem:schur-1-1} below,
$s_\lambda(q/u)=0$ unless $\lambda$ is a hook). By Lemma~\ref{lem:P11hooks},
$P(1,1,n)$ is exactly the set of hook partitions of length $\le n$, so
\eqref{eq:coeffs-Hilbert} captures precisely the hook coefficients we seek.
\end{proof}

Thus determining the hook coefficients $c_{(a,1^b)}(n)$ for all $m$ is equivalent
to extracting them from \eqref{eq:coeffs-Hilbert}.

\section{The hook super Schur function}\label{sec:superschur}

\begin{lemma}\label{lem:schur-1-1}
Let $\lambda$ be a partition.
\begin{enumerate}[\rm(i)]
\item $s_{\varnothing}(q/u)=1$.
\item If $\lambda=(a,1^b)$ with $a\ge1$ and $b\ge0$, then
$s_\lambda(q/u)= q^a u^b + q^{a-1}u^{b+1}$.
\item If $\lambda$ is not a hook (equivalently $\lambda_2\ge2$), then
$s_\lambda(q/u)=0$.
\end{enumerate}
\end{lemma}

\begin{proof}
We evaluate \eqref{eq:superschur-def}, recalling two single-variable facts.

\emph{Fact A.} In one commuting variable $q$, $s_\nu(q)=0$ unless $\nu$ is a
single row $\nu=(r)$ (including $r=0$, i.e.\ $\nu=\varnothing$), in which case
$s_{(r)}(q)=q^{r}$. Indeed $s_\nu(q)$ is the generating function of column-strict
fillings of $\nu$ with entries in $\{1\}$; such a filling exists iff every column
has length $\le1$, i.e.\ $\nu$ has a single row, and then there is exactly one
filling, of content $r$.

\emph{Fact B.} In one (commuting) variable $u$, the skew Schur function
$s_{\mu/\nu}(u)$ equals $u^{|\mu/\nu|}$ if $\mu/\nu$ is a horizontal strip
($\nu_i\le\mu_i$ and $\mu_{i+1}\le\nu_i$ for all $i$), and is $0$ otherwise; this
is the count of semistandard fillings of $\mu/\nu$ with entries in $\{1\}$, which
is $1$ exactly when $\mu/\nu$ is a horizontal strip.

In \eqref{eq:superschur-def} the second factor is $s_{\lambda'/\nu'}(u)$, a skew
Schur in the single variable $u$. By Fact B it is nonzero iff $\lambda'/\nu'$ is a
horizontal strip, equivalently $\lambda/\nu$ is a \emph{vertical} strip (conjugate
of a horizontal strip). Combining with Fact A: a term survives in
\eqref{eq:superschur-def} iff
\begin{equation}\label{eq:survival}
\nu=(r)\ \text{is a single row, \quad and\quad } \lambda/\nu \text{ is a vertical strip.}
\end{equation}

\emph{(i)} For $\lambda=\varnothing$ the only $\nu$ is $\varnothing=(0)$, with
$s_\varnothing(q)=1$ and $s_{\varnothing/\varnothing}(u)=1$, giving
$s_\varnothing(q/u)=1$.

\emph{(ii)} Let $\lambda=(a,1^b)$, $a\ge1$. A single-row $\nu=(r)$ satisfies
$\nu\subseteq\lambda$ iff $0\le r\le a$. The skew shape $\lambda/(r)$ has
$a-r$ cells in row $1$ and one cell in each of rows $2,\dots,b+1$. This is a
vertical strip (at most one cell per row) iff $a-r\le1$, i.e.\ $r\in\{a-1,a\}$
(and $r=a-1$ requires $a\ge1$, which holds).
\begin{itemize}
\item $r=a$: $\nu=(a)$, $|\lambda/\nu|=b$, contributing
$q^{a}\cdot u^{b}$.
\item $r=a-1$: $\nu=(a-1)$, $|\lambda/\nu|=b+1$, contributing
$q^{a-1}\cdot u^{b+1}$.
\end{itemize}
For each surviving $\nu$ the corresponding $\lambda/\nu$ is indeed a vertical
strip, so its single-variable skew Schur is $u^{|\lambda/\nu|}$ by Fact B. Hence
$s_{(a,1^b)}(q/u)=q^a u^b+q^{a-1}u^{b+1}$.

\emph{(iii)} Suppose $\lambda$ is not a hook, so $\lambda_2\ge2$. For any
single-row $\nu=(r)$ we have $\nu_i=0$ for $i\ge2$, while $\lambda$ has
$\lambda_2\ge2$ cells in its second row; hence row $2$ of $\lambda/\nu$ contains
$\lambda_2\ge2$ cells, so $\lambda/\nu$ is not a vertical strip. By
\eqref{eq:survival} no term survives and $s_\lambda(q/u)=0$.
\end{proof}

\section{The alternating formula for the hook coefficients}\label{sec:alternating}

We will use the following enumeration of Sagan--Swanson.

\begin{proposition}[\!\!{\cite[Theorem 5.1]{SaganSwanson2024}}]\label{prop:OSP-inv}
For $n,k\ge0$,
\begin{equation}\label{eq:OSP-equation}
[k]_q!\,S[n,k] = \sum_{w \in \OSP(n,k)}q^{\inv(w)},
\qquad\text{i.e.}\qquad
\langle q^i\rangle\,[k]_q!\,S[n,k]=\#\{w\in\OSP(n,k):\inv(w)=i\}.
\end{equation}
\end{proposition}

We also use the Rhoades--Wilson Hilbert series.

\begin{theorem}[\!\!\cite{RhoadesWilson2023}]\label{thm:hilbert}
For $n\ge1$,
\begin{equation}\label{eq:hilbert-rhoades-wilson}
\Hilb(R_n^{(1,1)};q;u) = \sum_{k=0}^n [k]_q!\,S[n,k]\,u^{\,n-k}.
\end{equation}
\end{theorem}

\begin{lemma}\label{lem:c_{(a,1^b)}}
For $a\ge1$ and $b\ge0$,
\begin{equation}\label{eq:first-formula-coeffs}
c_{(a,1^b)}(n) =
\sum_{i=0}^b (-1)^i\,\#\{w \in \OSP(n,n-b+i)\ :\ \inv(w) = a+i \}.
\end{equation}
\end{lemma}

\begin{proof}
Equating \eqref{eq:coeffs-Hilbert} and \eqref{eq:hilbert-rhoades-wilson} and
applying Lemma~\ref{lem:schur-1-1} to the left-hand side (only hooks and
$\varnothing$ contribute, by Lemma~\ref{lem:P11hooks}), we get
\begin{equation}\label{eq:master}
1 + \sum_{\substack{a \geq 1\\ b \geq 0}} c_{(a,1^b)}(n)\,
\big(q^a u^b+q^{a-1}u^{b+1}\big)
= \sum_{d=0}^n [d]_q!\,S[n,d]\,u^{n-d}.
\end{equation}
By \eqref{eq:OSP-equation} the right-hand side equals
\begin{equation}\label{eq:rhs-osp}
\sum_{d=0}^n \sum_{i\ge0} \#\{w\in\OSP(n,d):\inv(w)=i\}\,q^i\,u^{n-d}.
\end{equation}

We now extract the coefficient of $q^a u^b$ on both sides
($a\ge1$, $b\ge0$). On the left, a monomial $q^au^b$ arises from exactly two
hook terms: the $q^{a}u^{b}$ term of $(a,1^b)$ (present since $a\ge1$, $b\ge0$),
and, when $b\ge1$, the $q^{(a+1)-1}u^{(b-1)+1}=q^au^b$ term of
$(a+1,1^{b-1})$. No other hook contributes: the $(\alpha,1^\beta)$ term
$q^{\alpha}u^{\beta}+q^{\alpha-1}u^{\beta+1}$ equals $q^au^b$ only if
$(\alpha,\beta)=(a,b)$ (first monomial) or $(\alpha,\beta)=(a+1,b-1)$ (second
monomial); the constant $1$ contributes only to $q^0u^0$ and is irrelevant since
$a\ge1$. On the right, the coefficient of $q^au^b$ in \eqref{eq:rhs-osp} comes
from the stratum $d=n-b$ (so that $u^{n-d}=u^{b}$) and $i=a$.

Hence, for $b\ge1$,
\begin{equation}\label{eq:recurrence}
c_{(a,1^b)}(n)+c_{(a+1,1^{b-1})}(n)
= \#\{w\in\OSP(n,n-b):\inv(w)=a\},
\end{equation}
and, for $b=0$,
\begin{equation}\label{eq:base}
c_{(a)}(n) = \#\{w\in\OSP(n,n):\inv(w)=a\}.
\end{equation}

We solve the recurrence by telescoping in the index $b$. Define, for $i\ge0$,
\begin{equation}
A_i := \#\{w\in\OSP(n,n-b+i):\inv(w)=a+i\}.
\end{equation}
We claim $c_{(a,1^b)}(n)=\sum_{i=0}^b(-1)^iA_i$. Apply
\eqref{eq:recurrence} with the substitutions $(a,b)\mapsto(a+i,b-i)$, valid for
$0\le i\le b-1$ since then $b-i\ge1$:
\begin{equation}\label{eq:tele}
c_{(a+i,1^{b-i})}(n)
= \#\{w\in\OSP(n,n-(b-i)):\inv(w)=a+i\}
- c_{(a+i+1,1^{b-i-1})}(n)
= A_i - c_{(a+i+1,1^{b-i-1})}(n).
\end{equation}
Starting from $i=0$ and iterating \eqref{eq:tele} for $i=0,1,\dots,b-1$,
\begin{align*}
c_{(a,1^b)}(n)
&= A_0 - c_{(a+1,1^{b-1})}(n)
= A_0 - A_1 + c_{(a+2,1^{b-2})}(n)
= \cdots\\
&= \sum_{i=0}^{b-1}(-1)^i A_i + (-1)^b c_{(a+b,1^0)}(n).
\end{align*}
By the base case \eqref{eq:base},
$c_{(a+b)}(n)=\#\{w\in\OSP(n,n):\inv(w)=a+b\}=A_b$ (taking $i=b$ above:
$\OSP(n,n-b+b)=\OSP(n,n)$ and $\inv=a+b$). Therefore
$c_{(a,1^b)}(n)=\sum_{i=0}^b(-1)^iA_i$, which is
\eqref{eq:first-formula-coeffs}.
\end{proof}

\section{The Sagan--Swanson involution and a manifestly positive formula}
\label{sec:phi}

\subsection{The involution $\phi$}
We recall the involution $\phi$ of Sagan--Swanson
\cite[Section 5.1]{SaganSwanson2024}. Given $w=(S_1/\cdots/S_k)\in\OSP(n,k)$ and
$i\in\{1,\dots,k\}$, write $M=\max S_i$. We say $w$ is:
\begin{itemize}
\item \textbf{splittable} at $M=\max S_i$ if $|S_i|>1$;
\item \textbf{mergeable} at $M=\max S_i$ (for $i\le k-1$) if $|S_i|=1$ and
$M>\max S_{i+1}$.
\end{itemize}
Find the largest $M$ at which $w$ is splittable or mergeable, if one exists. Then
$w$ is of \textbf{type I} if it is mergeable at this $M$; \textbf{type II} if it
is splittable at this $M$; \textbf{type III} if no such $M$ exists. Define
\begin{equation}\label{eq:phi}
\phi(w)=\begin{cases}
(S_1/\cdots/S_{i-1}/S_i\cup S_{i+1}/S_{i+2}/\cdots/S_k)
& \text{type I (merge at $S_i$)},\\[2pt]
(S_1/\cdots/S_{i-1}/\{M\}/S_i\setminus\{M\}/S_{i+1}/\cdots/S_k)
& \text{type II (split $S_i$ at $M$)},\\[2pt]
w & \text{type III}.
\end{cases}
\end{equation}

\begin{lemma}\label{lem:phi-props}
$\phi$ is an involution on $\bigcup_{k}\OSP(n,k)$. If $w$ is of type I then
$\phi(w)$ is of type II with $\blocks(\phi(w))=\blocks(w)-1$ and
$\inv(\phi(w))=\inv(w)-1$; if $w$ is of type II then $\phi(w)$ is of type I with
$\blocks(\phi(w))=\blocks(w)+1$ and $\inv(\phi(w))=\inv(w)+1$. Type III elements
are the fixed points. Consequently $n-\blocks(w)-\inv(w)$ and the parity of
$\blocks(w)$ behave as stated, and the quantity $\inv(w)+\blocks(w)$ increases by
$2$ under a split and decreases by $2$ under a merge.
\end{lemma}

\begin{proof}
That $\phi$ is a well-defined involution exchanging types I and II at the same
value $M$, with the stated effect on $\blocks$ and $\inv$, is
\cite[Section 5.1]{SaganSwanson2024}. The merge in type I deletes the singleton
block $\{M\}$ and absorbs it into the next block; this removes one block and,
since $\{M\}$ contributed exactly one inversion pair (the merge condition
$M>\max S_{i+1}$ means $M$ created precisely the inversion with the following
block and none thereafter, by maximality of $M$), removes one inversion. The
split is the inverse operation, adding one block and one inversion. Hence
$\blocks$ and $\inv$ change by $\mp1$ together, and $\phi$ is its own inverse. The
remaining bookkeeping statements are immediate.
\end{proof}

\subsection{The signed set and its cancellation}
Fix integers $\ell\ge1$ and $0\le b\le\lfloor(\ell-1)/2\rfloor$. Define
\begin{equation}\label{eq:Sdef}
S_{n,\ell,b}=\bigcup_{c=0}^b \{w\in\OSP(n,n-c):\inv(w)=\ell-b-c\},
\end{equation}
a disjoint union over $c$ (the strata have distinct block counts $n-c$). On
$S_{n,\ell,b}$ we use the sign
\begin{equation}\label{eq:epsdef}
\varepsilon(w):=(-1)^{\,\blocks(w)-(n-b)}=(-1)^{\,b-c}\qquad
\text{for }w\in\OSP(n,n-c).
\end{equation}

\begin{lemma}[Standing window]\label{lem:window}
If $w\in S_{n,\ell,b}$ then $n-b\le\blocks(w)\le n$ and
$\ell-2b\le\inv(w)\le\ell-b$.
\end{lemma}

\begin{proof}
By \eqref{eq:Sdef}, $w\in\OSP(n,n-c)$ with $0\le c\le b$, so
$n-b\le\blocks(w)=n-c\le n$, and $\inv(w)=\ell-b-c\in[\ell-2b,\ell-b]$.
\end{proof}

\begin{lemma}[Classification of $\phi$ on $S_{n,\ell,b}$]\label{lem:classify}
Let $w\in S_{n,\ell,b}$, say $w\in\OSP(n,n-c)$ with $\inv(w)=\ell-b-c$,
$0\le c\le b$. Then:
\begin{enumerate}[\rm(i)]
\item if $w$ is type I and $c=b$ (equivalently $\inv(w)=\ell-2b$), then
$\phi(w)\notin S_{n,\ell,b}$;
\item if $w$ is type I and $c<b$, then $\phi(w)\in S_{n,\ell,b}$, lying in
$\OSP(n,n-(c+1))$ with $\inv(\phi(w))=\ell-b-(c+1)$;
\item if $w$ is type II, then $c>0$ and $\phi(w)\in S_{n,\ell,b}$, lying in
$\OSP(n,n-(c-1))$ with $\inv(\phi(w))=\ell-b-(c-1)$;
\item $w$ cannot be type III.
\end{enumerate}
In cases (ii) and (iii), $\varepsilon(\phi(w))=-\varepsilon(w)$.
\end{lemma}

\begin{proof}
\emph{(iv)} Suppose $w$ were type III, i.e.\ a fixed point of $\phi$. A type III
ordered set partition has every block a singleton arranged in increasing order of
maxima, hence has $n$ blocks and $\inv=0$; in particular $\blocks(w)=n$, forcing
$c=0$, and $\inv(w)=0$, forcing $\ell-b=0$, i.e.\ $\ell=b$. But $\ell\ge1$ gives
$b=\ell>\lfloor(\ell-1)/2\rfloor$, contradicting $b\le\lfloor(\ell-1)/2\rfloor$.
So $w$ is type I or II.

\emph{(i)--(ii) (type I).} By Lemma~\ref{lem:phi-props}, $\phi(w)$ is type II with
$\blocks(\phi(w))=\blocks(w)-1=n-(c+1)$ and $\inv(\phi(w))=\inv(w)-1=\ell-b-c-1$.
For $\phi(w)$ to lie in $S_{n,\ell,b}$ we need its block count $n-(c+1)$ to be a
valid stratum, i.e.\ $0\le c+1\le b$, and its inversion number to match that
stratum, i.e.\ $\inv(\phi(w))=\ell-b-(c+1)$. The inversion condition holds
identically: $\ell-b-c-1=\ell-b-(c+1)$. Thus membership is governed solely by
$c+1\le b$:
\begin{itemize}
\item if $c=b$, then $c+1=b+1>b$, so $\phi(w)\notin S_{n,\ell,b}$; this is
(i), and indeed $\inv(w)=\ell-b-b=\ell-2b$.
\item if $c<b$, then $0\le c+1\le b$, so $\phi(w)\in\OSP(n,n-(c+1))$ with
$\inv(\phi(w))=\ell-b-(c+1)$, hence $\phi(w)\in S_{n,\ell,b}$; this is (ii).
\end{itemize}

\emph{(iii) (type II).} By Lemma~\ref{lem:phi-props}, $\phi(w)$ is type I with
$\blocks(\phi(w))=\blocks(w)+1=n-(c-1)$ and $\inv(\phi(w))=\inv(w)+1=\ell-b-c+1$.
Since $w$ is splittable, it has a block of size $\ge2$, so $\blocks(w)=n-c<n$,
forcing $c\ge1$, i.e.\ $c>0$; in particular $0\le c-1\le b$. The inversion
condition for the stratum $\OSP(n,n-(c-1))$ requires
$\inv(\phi(w))=\ell-b-(c-1)$, which holds: $\ell-b-c+1=\ell-b-(c-1)$. Hence
$\phi(w)\in S_{n,\ell,b}$.

\emph{Sign.} In cases (ii) and (iii), $\blocks(\phi(w))=\blocks(w)\pm1$, so by
\eqref{eq:epsdef} $\varepsilon(\phi(w))=(-1)^{\pm1}\varepsilon(w)=-\varepsilon(w)$.
\end{proof}

\begin{lemma}[Survivors]\label{lem:survivors}
Let
\begin{equation}\label{eq:Tdef}
T_{n,\ell,b}=\{w\in\OSP(n,n-b)\ :\ w\text{ is type I and }\inv(w)=\ell-2b\}.
\end{equation}
Then $\phi$ restricts to a fixed-point-free, sign-reversing involution on
$S_{n,\ell,b}\setminus T_{n,\ell,b}$, and $\phi$ maps $T_{n,\ell,b}$ out of
$S_{n,\ell,b}$. Consequently
\begin{equation}\label{eq:signed-sum}
\sum_{w\in S_{n,\ell,b}}\varepsilon(w) = \#\,T_{n,\ell,b}.
\end{equation}
\end{lemma}

\begin{proof}
First, $T_{n,\ell,b}\subseteq S_{n,\ell,b}$: an element of $T_{n,\ell,b}$ lies in
$\OSP(n,n-b)$ (so $c=b$) with $\inv=\ell-2b=\ell-b-b$, matching \eqref{eq:Sdef}.
By Lemma~\ref{lem:classify}(i) every $w\in T_{n,\ell,b}$ (type I, $c=b$) has
$\phi(w)\notin S_{n,\ell,b}$.

Now take $w\in S_{n,\ell,b}\setminus T_{n,\ell,b}$. By
Lemma~\ref{lem:classify}(iv) $w$ is type I or II. If $w$ is type I, then $w\notin
T_{n,\ell,b}$ forces $c<b$ (since type-I elements with $c=b$ are exactly
$T_{n,\ell,b}$), so case (ii) applies; if $w$ is type II, case (iii) applies. In
both cases $\phi(w)\in S_{n,\ell,b}$ and $\varepsilon(\phi(w))=-\varepsilon(w)$.

We check $\phi(w)\in S_{n,\ell,b}\setminus T_{n,\ell,b}$ as well, so that $\phi$
genuinely restricts to this subset. In case (ii) ($w$ type I, $c<b$),
Lemma~\ref{lem:phi-props} gives $\phi(w)$ of type II, and type II elements are
never in $T_{n,\ell,b}$ (which consists of type I elements); so
$\phi(w)\in S_{n,\ell,b}\setminus T_{n,\ell,b}$. In case (iii) ($w$ type II,
$c>0$), $\phi(w)$ is type I in $\OSP(n,n-(c-1))$; it lies in $T_{n,\ell,b}$ only
if its block count is $n-b$, i.e.\ $c-1=b$, i.e.\ $c=b+1$, impossible since
$c\le b$; so again $\phi(w)\in S_{n,\ell,b}\setminus T_{n,\ell,b}$.

Since $\phi$ is an involution on all ordered set partitions
(Lemma~\ref{lem:phi-props}) and maps
$S_{n,\ell,b}\setminus T_{n,\ell,b}$ into itself, it restricts to an involution
there. It is fixed-point-free on this subset: a fixed point of $\phi$ is type III,
excluded by Lemma~\ref{lem:classify}(iv). As it reverses $\varepsilon$, the
contributions of $S_{n,\ell,b}\setminus T_{n,\ell,b}$ to the signed sum cancel in
pairs:
\begin{equation}
\sum_{w\in S_{n,\ell,b}}\varepsilon(w)
=\sum_{w\in T_{n,\ell,b}}\varepsilon(w)
+\sum_{w\in S_{n,\ell,b}\setminus T_{n,\ell,b}}\varepsilon(w)
=\sum_{w\in T_{n,\ell,b}}\varepsilon(w).
\end{equation}
Finally every $w\in T_{n,\ell,b}$ has $\blocks(w)=n-b$, so $c=b$ and
$\varepsilon(w)=(-1)^{(n-b)-(n-b)}=(-1)^0=+1$ by \eqref{eq:epsdef}. Hence
$\sum_{w\in S_{n,\ell,b}}\varepsilon(w)=\#\,T_{n,\ell,b}$.
\end{proof}

\subsection{The main theorem}

\begin{theorem}\label{thm:improved-coefficient-formula}
For $a\ge1$ and $b\ge0$, the hook coefficient is given, manifestly positively, by
\begin{equation}\label{eq:second-formula-coeffs}
c_{(a,1^b)}(n)=\#\{w\in\OSP(n,n-b)\ :\ \inv(w)=a,\ w\text{ is type I}\}.
\end{equation}
Moreover $c_\varnothing(n)=1$ and $c_\lambda(n)=0$ for every non-hook $\lambda$.
\end{theorem}

\begin{proof}
Set $\ell:=a+2b$. Then $\ell\ge1$ (as $a\ge1$) and
$b\le\lfloor(\ell-1)/2\rfloor$, because $2b=\ell-a\le\ell-1$ gives
$b\le(\ell-1)/2$ and $b$ is an integer. With this $\ell$, the defining set
\eqref{eq:Sdef} becomes
\begin{equation}
S_{n,a+2b,b}=\bigcup_{c=0}^b\{w\in\OSP(n,n-c):\inv(w)=a+b-c\}.
\end{equation}
Reindex the alternating formula \eqref{eq:first-formula-coeffs} by
$c:=b-i$ (so $i=b-c$ runs over $0,\dots,b$ as $c$ runs over $b,\dots,0$). Then
$\OSP(n,n-b+i)=\OSP(n,n-c)$, $\inv=a+i=a+b-c$, and the sign is
$(-1)^i=(-1)^{b-c}=\varepsilon(w)$ for $w\in\OSP(n,n-c)$, by \eqref{eq:epsdef}.
Therefore
\begin{equation}
c_{(a,1^b)}(n)
=\sum_{c=0}^b(-1)^{b-c}\#\{w\in\OSP(n,n-c):\inv(w)=a+b-c\}
=\sum_{w\in S_{n,a+2b,b}}\varepsilon(w).
\end{equation}
By Lemma~\ref{lem:survivors} with $\ell=a+2b$,
\begin{equation}
c_{(a,1^b)}(n)=\#\,T_{n,a+2b,b}
=\#\{w\in\OSP(n,n-b):w\text{ type I},\ \inv(w)=(a+2b)-2b\},
\end{equation}
and $(a+2b)-2b=a$, which is exactly \eqref{eq:second-formula-coeffs}. This count
is a cardinality, hence a nonnegative integer, so the formula is manifestly
positive.

The statement $c_\varnothing(n)=1$ is the coefficient of $q^0u^0$ in
\eqref{eq:master}: the constant term on the left is $1$ and on the right is the
$d=n$, $i=0$ contribution $\#\{w\in\OSP(n,n):\inv(w)=0\}=1$ (the unique increasing
ordered set partition into singletons), confirming $c_\varnothing(n)=1$. For a
non-hook $\lambda$ ($\lambda_2\ge2$) we have $\lambda\notin P(1,1,n)$, so it does
not occur in \eqref{eq:coeffs-Hilbert}; equivalently $s_\lambda(q/u)=0$ by
Lemma~\ref{lem:schur-1-1}(iii). By the universality in
Proposition~\ref{prop:bergeron} and the supersymmetric identification in
Proposition~\ref{prop:ds}, the coefficient $c_\lambda(n)$ of every non-hook
$\lambda$ is $0$.
\end{proof}

\section{Example: $n=4$}
The following table lists the nonzero hook coefficients $c_\lambda(4)$ together
with the ordered set partitions counted by
Theorem~\ref{thm:improved-coefficient-formula}. (All blocks are written
increasingly; e.g.\ $2|1|3|4$ is $(\{2\}/\{1\}/\{3\}/\{4\})$.)

\begin{figure}[h]
\begin{tabular}{l|l|llllll}
\text{Coefficient} & \text{Value} & \text{OSPs}         \\ \hline
$c_{(1)}(4)$       & 3            & $2|1|3|4$  & $1|3|2|4$ & $1|2|4|3$ &           &           &           \\
$c_{(1,1)}(4)$     & 3            & $1|4|23$   & $12|4|3$  & $3|12|4$  &           &           &           \\
$c_{(1,1,1)}(4)$   & 1            & $4|123$    &           &           &           &           &           \\
$c_{(2)}(4)$       & 5            & $2|1|4|3$  & $1|3|4|2$ & $2|3|1|4$ & $3|1|2|4$ & $1|4|2|3$ &           \\
$c_{(2,1)}(4)$     & 5            & $4|1|23$   & $13|4|2$  & $2|4|13$  & $3|4|12$  & $4|12|3$  &           \\
$c_{(3)}(4)$       & 6            & $2|3|4|1$  & $3|2|1|4$ & $2|4|1|3$ & $3|1|4|2$ & $4|1|2|3$ & $1|4|3|2$ \\
$c_{(3,1)}(4)$     & 4            & $4|13|2$   & $23|4|1$  & $4|2|13$  & $4|3|12$  &           &           \\
$c_{(4)}(4)$       & 5            & $4|1|3|2$  & $4|2|1|3$ & $3|4|1|2$ & $3|2|4|1$ & $2|4|3|1$ &           \\
$c_{(4,1)}(4)$     & 1            & $4|23|1$   &           &           &           &           &           \\
$c_{(5)}(4)$       & 3            & $4|2|3|1$  & $4|3|1|2$ & $3|4|2|1$ &           &           &           \\
$c_{(6)}(4)$       & 1            & $4|3|2|1$  &           &           &           &           &          
\end{tabular}
\caption{The nonzero hook coefficients $c_\lambda(4)$ and the ordered set
partitions counted by Theorem~\ref{thm:improved-coefficient-formula}. Apart from
$c_\varnothing(4)=1$, all other $c_\lambda(4)$ are $0$.}
\end{figure}

These reproduce the Schur (super Schur) expansion
\begin{align*}
\Hilb(R_4^{(1,1)};q;u)
&= (q^6+3q^5+5q^4+6q^3+5q^2+3q+1)
 + (q^5+4q^4+9q^3+11q^2+8q+3)u\\
&\quad + (q^3+4q^2+6q+3)u^2 + u^3\\
&= 1+\sum_{a=1}^{6} c_{(a)}(4)\,s_{(a)}(q/u)
 + \sum_{a=1}^{4} c_{(a,1)}(4)\,s_{(a,1)}(q/u)
 + c_{(1,1,1)}(4)\,s_{(1,1,1)}(q/u),
\end{align*}
using $s_{(a,1^b)}(q/u)=q^au^b+q^{a-1}u^{b+1}$ from
Lemma~\ref{lem:schur-1-1}.

\begin{thebibliography}{99}
\bibitem{Bergeron2020} Fran\c{c}ois Bergeron, \emph{The bosonic-fermionic
diagonal coinvariant modules conjecture}, preprint (2020),
\url{https://arxiv.org/abs/2005.00924}.
\bibitem{ChengWangBook} Shun-Jen Cheng and Weiqiang Wang, \emph{Dualities and
representations of Lie superalgebras}, Graduate Studies in Mathematics, vol.\
144, American Mathematical Society, Providence, RI, 2012.
\bibitem{Lentfer-Supersymmetry} John Lentfer, \emph{Diagonal supersymmetry for
coinvariant rings}, preprint (2025), \url{https://arxiv.org/abs/2505.14885}.
\bibitem{RhoadesWilson2023} Brendon Rhoades and Andrew Timothy Wilson, \emph{The
Hilbert series of the superspace coinvariant ring}, Forum Math.\ Pi 12 (2024),
Id/No e16.
\bibitem{SaganSwanson2024} Bruce E. Sagan and Joshua P. Swanson, \emph{$q$-Stirling
numbers in type $B$}, Eur.\ J.\ Comb.\ 118 (2024), Id/No 103899.
\end{thebibliography}

\end{document}
